% =====================================================================
%  A parsimonious breakpoint panel for meningioma recurrence
%
%  Compile:  pdflatex breakpoint_panel.tex   (twice, for references)
%  Figures are expected in ../figures_nejm/ relative to this file.
%
%  Every numeric value below was read from the generated CSVs in
%  ../ ; see \S\ref{sec:repro} for the exact provenance of each table.
% =====================================================================
\documentclass[11pt]{article}

\usepackage[margin=1in]{geometry}
\usepackage{graphicx}
\usepackage{booktabs}
\usepackage{caption}
\usepackage{subcaption}
\usepackage{siunitx}
\usepackage{xcolor}
\usepackage{amsmath}
\usepackage[colorlinks=true,linkcolor=black,citecolor=black,urlcolor=blue]{hyperref}
\usepackage{float}

\graphicspath{{../figures_nejm/}{./figures_nejm/}{./}}

% Let \bfseries take effect inside siunitx S columns, and keep the parser from
% choking on the bracketed difference values in Table 3.
\sisetup{detect-weight=true, detect-family=true, table-align-text-before=false}

\captionsetup{font=small,labelfont=bf,justification=justified,singlelinecheck=false}
\setlength{\parskip}{0.4em}
\setlength{\parindent}{0pt}

\definecolor{nejmred}{HTML}{BC3C29}
\newcommand{\note}[1]{\textcolor{gray}{\small #1}}

\title{\vspace{-2em}A parsimonious breakpoint panel for meningioma recurrence:\\
\large eight genomic windows compared with WHO and integrated grade}
\author{Breakpoint Architecture Analysis}
\date{4 August 2026}

\begin{document}
\maketitle

\begin{abstract}
\noindent
Eight binary breakpoint indicators, each defined by a fixed genomic interval and
discovered entirely within training folds, were compared with WHO grade,
integrated grade and arm-level copy-number alteration (SCNA) for prediction of
meningioma recurrence at 3, 5 and 7 years. The panel alone matched integrated
grade at 5 years (AUC 0.735 vs 0.735) and exceeded it at 7 years
(0.752 vs 0.733), while every conventional marker declined across the same
interval. Added to arm-level SCNA the panel improved discrimination by
+0.020 to +0.030 AUC in the broad cohort and by +0.041 to +0.081 in a restricted
surgical cohort. Combined with five routine clinical covariates it reached
AUC 0.850 at 3 years using 60 features, and 0.836 using 13. The panel
independently recovered the known chromosome 1p 112--125\,Mb hotspot in
88--100\% of training folds. It did \emph{not} replicate on an external
methylation-derived cohort, for a measured and specific reason: assay
resolution.
\end{abstract}

% =====================================================================
\section{Cohort and methods}
% =====================================================================

Institutional meningioma cohort, GRCh37, event-called segmental copy number.
Recurrence was evaluated at fixed horizons under a strict endpoint: recurrence
on or before the horizon counted as an event, event-free follow-up through the
horizon as a non-event, and censoring before the horizon without recurrence led
to exclusion. Cohort sizes are given in Table~\ref{tab:cohort}.

\begin{table}[H]
\centering
\caption{Analytic cohorts. The restricted subset requires primary presentation,
gross-total resection, no prior radiotherapy, no adjuvant radiotherapy, not
radiation-induced, and non-angiomatous histology --- all six conditions.}
\label{tab:cohort}
\begin{tabular}{lrrrr}
\toprule
& \multicolumn{2}{c}{Broad cohort} & \multicolumn{2}{c}{Restricted cohort} \\
\cmidrule(lr){2-3}\cmidrule(lr){4-5}
Horizon & Tumours & Recurrences & Tumours & Recurrences \\
\midrule
3 years & 904 & 204 & 387 & 22 \\
5 years & 674 & 264 & 248 & 37 \\
7 years & 514 & 296 & 155 & 43 \\
\bottomrule
\end{tabular}
\end{table}

\paragraph{Breakpoint windows.}
Each chromosome arm was divided into 20 bins of equal arm-relative width
(0 at the telomere, 1 at the centromere). A window is scored 1 for a tumour if
any breakpoint falls inside it. Windows present in fewer than 20 tumours were
discarded, leaving 745 candidates.

\paragraph{Window selection is internal to training folds.}
Within every training fold, candidate windows were ranked by a
$\chi^2$ test of association with the outcome and the top $k=8$ retained. The
model was then fitted on those folds only and scored on held-out patients.
Reported AUCs therefore include the cost of discovery. No frozen or previously
published window list was used as an input.

\paragraph{Validation.}
Five repeats of patient-grouped 5-fold cross-validation. All tumours from one
patient are assigned to the same fold, which matters because 75 patients
contribute more than one resection. Models are binomial elastic net
($\alpha=0.25$), predictors standardised and missing clinical values imputed
using training-fold statistics only.

\paragraph{Restricted-cohort estimates use transport, not refitting.}
Following the locked project's design, the model is fitted on the broad cohort
with patient-grouped out-of-fold prediction and those honest scores are then
subset to restricted tumours. Refitting inside the restricted subset was
attempted and abandoned: with 22 events at 3 years it produced an arm-level SCNA
AUC of 0.500, which is a statement about sample size rather than about biology.

% =====================================================================
\section{Discrimination}
% =====================================================================

\begin{figure}[H]
\centering
\includegraphics[width=\textwidth]{figure1_discrimination.png}
\caption{Discrimination of the eight-window panel against WHO grade and
integrated grade. Comparator AUCs are computed \emph{separately in each cohort};
published values are broad-cohort figures and differ materially in the
restricted subset. Panel labels give tumours / recurrences.}
\label{fig:disc}
\end{figure}

\begin{table}[H]
\centering
\caption{Area under the ROC curve. The panel alone uses eight binary indicators
and no other information.}
\label{tab:disc}
\sisetup{table-format=1.3}
\begin{tabular}{llSSSS}
\toprule
Cohort & Horizon & {WHO grade} & {Integrated grade} & {Arm-level SCNA} & {Panel alone} \\
\midrule
Broad      & 3 years & 0.690 & 0.766 & 0.777 & 0.702 \\
           & 5 years & 0.656 & 0.735 & 0.766 & 0.735 \\
           & 7 years & 0.635 & 0.733 & 0.767 & 0.752 \\
\addlinespace
Restricted & 3 years & 0.686 & 0.827 & 0.796 & 0.715 \\
           & 5 years & 0.647 & 0.782 & 0.766 & 0.753 \\
           & 7 years & 0.622 & 0.752 & 0.740 & 0.757 \\
\bottomrule
\end{tabular}
\end{table}

Two features of Table~\ref{tab:disc} are worth separating.

First, \textbf{the panel alone is comparable to integrated grade, not superior
to arm-level SCNA}. At 5 years the two are identical (0.735 vs 0.735); at
7 years the panel is ahead in both cohorts (0.752 vs 0.733 broad; 0.757 vs 0.752
restricted). It does not beat arm-level SCNA at 3 or 5 years.

Second, \textbf{the markers move in opposite directions with horizon}
(Table~\ref{tab:trend}). Every conventional marker loses discrimination as the
horizon lengthens; the breakpoint panel gains it. This is the clearest evidence
in the analysis that breakpoint architecture indexes a different prognostic axis
from grade.

\begin{table}[H]
\centering
\caption{Change in AUC from 3 to 7 years, both cohorts. Every conventional
marker declines; the breakpoint panel is the only one that improves.}
\label{tab:trend}
\begin{tabular}{lcccccc}
\toprule
& \multicolumn{3}{c}{Broad cohort} & \multicolumn{3}{c}{Restricted cohort} \\
\cmidrule(lr){2-4}\cmidrule(lr){5-7}
Marker & 3 yr & 7 yr & Change & 3 yr & 7 yr & Change \\
\midrule
WHO grade                 & 0.690 & 0.635 & $-0.055$ & 0.686 & 0.622 & $-0.064$ \\
Integrated grade          & 0.766 & 0.733 & $-0.033$ & 0.827 & 0.752 & $-0.075$ \\
Arm-level SCNA            & 0.777 & 0.767 & $-0.010$ & 0.796 & 0.740 & $-0.056$ \\
\textbf{Breakpoint panel} & 0.702 & 0.752 & $\mathbf{+0.050}$ & 0.715 & 0.757 & $\mathbf{+0.042}$ \\
\bottomrule
\end{tabular}
\end{table}

The divergence is present in both cohorts and is if anything larger in the
restricted one, where integrated grade loses 0.075 across the interval while the
panel gains 0.042.

% =====================================================================
\section{Incremental value over arm-level SCNA}
% =====================================================================

\begin{figure}[H]
\centering
\includegraphics[width=\textwidth]{figure1b_discrimination_with_combined.png}
\caption{As Figure~\ref{fig:disc}, with the combined SCNA + panel model added as
a fourth bar, so the panel can be read both on its own and as part of a full
genomic model.}
\label{fig:disc1b}
\end{figure}

\begin{figure}[H]
\centering
\includegraphics[width=\textwidth]{figure3_scna_plus_breakpoints.png}
\caption{Arm-level SCNA with and without the eight-window panel. Top row:
discrimination. Bottom row: specificity at a matched sensitivity of
approximately 80\%. Left, broad cohort; right, restricted cohort.}
\label{fig:incr}
\end{figure}

\begin{table}[H]
\centering
\caption{Incremental discrimination from adding eight windows to arm-level SCNA.}
\label{tab:incr}
\sisetup{table-format=1.3}
\begin{tabular}{llSSS}
\toprule
Cohort & Horizon & {SCNA} & {SCNA + panel} & {Difference} \\
\midrule
Broad      & 3 years & 0.777 & 0.780 & {+0.003} \\
           & 5 years & 0.766 & 0.797 & {+0.030} \\
           & 7 years & 0.767 & 0.797 & {+0.030} \\
\addlinespace
Restricted & 3 years & 0.796 & 0.804 & {+0.007} \\
           & 5 years & 0.766 & 0.847 & {+0.081} \\
           & 7 years & 0.740 & 0.812 & {+0.072} \\
\bottomrule
\end{tabular}
\end{table}

The increment is roughly \textbf{2.5 times larger in the restricted cohort} at
5 and 7 years. A plausible reading is that recurrence in the broad cohort is
driven substantially by treatment and presentation --- subtotal resection, prior
radiotherapy, radiation-induced tumours --- and that removing those leaves
variation closer to intrinsic tumour biology. The 3-year increment is negligible
in both cohorts.

% =====================================================================
\section{Operating characteristics}
% =====================================================================

\begin{figure}[H]
\centering
\includegraphics[width=\textwidth]{figure2_specificity.png}
\caption{Specificity at a matched sensitivity of approximately 80\%. WHO grade
and integrated grade are three-level ordinals: the lowest threshold that reaches
80\% sensitivity classifies essentially every tumour as high risk, leaving no
specificity. Achieved sensitivity for both grades is 1.00.}
\label{fig:spec}
\end{figure}

\begin{table}[H]
\centering
\caption{Specificity at matched sensitivity ($\approx$80\%). Grade values are
exactly zero because 80\% sensitivity is unreachable without flagging the whole
cohort.}
\label{tab:spec}
\sisetup{table-format=1.3}
\begin{tabular}{llSSSS}
\toprule
Cohort & Horizon & {WHO} & {Integrated grade} & {SCNA} & {SCNA + panel} \\
\midrule
Broad      & 3 years & 0.000 & 0.000 & 0.636 & 0.660 \\
           & 5 years & 0.000 & 0.000 & 0.473 & 0.610 \\
           & 7 years & 0.000 & 0.000 & 0.505 & 0.642 \\
\addlinespace
Restricted & 3 years & 0.000 & 0.000 & 0.567 & 0.778 \\
           & 5 years & 0.000 & 0.000 & 0.450 & 0.787 \\
           & 7 years & 0.000 & 0.000 & 0.429 & 0.679 \\
\bottomrule
\end{tabular}
\end{table}

This is a granularity limitation of grading rather than an artefact of the
analysis, and it has a direct clinical reading: a three-level ordinal cannot be
operated at a chosen sensitivity, whereas a continuous score can be thresholded
anywhere. In the restricted cohort the combined model raises specificity from
0.450 to 0.787 at 5 years at the same sensitivity.

% =====================================================================
\section{Reclassification of grade-discordant tumours}
% =====================================================================

The clinically pointed question is not overall accuracy but what happens to the
tumours grading gets wrong. Two discordant groups were defined at each horizon:
tumours \emph{over-called} by grade (grade 2/3, never recurred), which a useful
model should place below its high-risk threshold; and tumours
\emph{under-called} by grade (grade 1, recurred), which it should place above.
All models were thresholded at approximately 80\% sensitivity.

\begin{figure}[H]
\centering
\includegraphics[width=\textwidth]{figure5_discordance_broad.png}
\caption{Reclassification of grade-discordant tumours, broad cohort.
Panels A--B, versus WHO grade; C--D, versus integrated grade. Left column,
over-called by grade; right column, under-called. $n$ is the number of
discordant tumours in each cell.}
\label{fig:disc2}
\end{figure}

\begin{figure}[H]
\centering
\includegraphics[width=\textwidth]{figure5_discordance_restricted.png}
\caption{Reclassification of grade-discordant tumours, restricted cohort.
Layout as Figure~\ref{fig:disc2}. Cell sizes are small (5--73 tumours) and the
bars should be read accordingly.}
\label{fig:disc2r}
\end{figure}

\begin{table}[H]
\centering
\caption{Fraction of grade-discordant tumours correctly reclassified, versus
\textbf{integrated grade}, both cohorts.}
\label{tab:discord}
\sisetup{table-format=1.3}
\begin{tabular}{llrSSrSS}
\toprule
& & \multicolumn{3}{c}{Broad cohort} & \multicolumn{3}{c}{Restricted cohort} \\
\cmidrule(lr){3-5}\cmidrule(lr){6-8}
Direction & Horizon & {$n$} & {SCNA} & {Panel} & {$n$} & {SCNA} & {Panel} \\
\midrule
Over-called    & 3 years & 180 & 0.044 & \bfseries 0.156 & 62 & 0.000 & \bfseries 0.161 \\
(grade 2/3,    & 5 years & 107 & 0.047 & \bfseries 0.196 & 34 & 0.000 & \bfseries 0.235 \\
never recurred)& 7 years & 52  & 0.096 & \bfseries 0.250 & 19 & 0.000 & \bfseries 0.316 \\
\addlinespace
Under-called   & 3 years & 55  & 0.309 & \bfseries 0.473 & 5  & 0.200 & \bfseries 0.400 \\
(grade 1,      & 5 years & 87  & 0.448 & \bfseries 0.563 & 12 & 0.500 & \bfseries 0.667 \\
recurred)      & 7 years & 102 & 0.471 & \bfseries 0.539 & 16 & 0.562 & \bfseries 0.688 \\
\bottomrule
\end{tabular}
\end{table}

\begin{table}[H]
\centering
\caption{The same comparison versus \textbf{WHO grade}. The three models are
close to indistinguishable, in contrast to Table~\ref{tab:discord}.}
\label{tab:discord_who}
\sisetup{table-format=1.3}
\begin{tabular}{llrSSrSS}
\toprule
& & \multicolumn{3}{c}{Broad cohort} & \multicolumn{3}{c}{Restricted cohort} \\
\cmidrule(lr){3-5}\cmidrule(lr){6-8}
Direction & Horizon & {$n$} & {SCNA} & {Panel} & {$n$} & {SCNA} & {Panel} \\
\midrule
Over-called    & 3 years & 196 & 0.347 & 0.362 & 73 & 0.452 & 0.452 \\
(grade 2/3,    & 5 years & 117 & 0.333 & 0.342 & 38 & 0.368 & 0.447 \\
never recurred)& 7 years & 66  & 0.455 & 0.424 & 21 & 0.476 & 0.524 \\
\addlinespace
Under-called   & 3 years & 75  & 0.573 & 0.640 & 10 & 0.700 & 0.800 \\
(grade 1,      & 5 years & 113 & 0.646 & 0.673 & 20 & 0.750 & 0.750 \\
recurred)      & 7 years & 134 & 0.634 & 0.679 & 25 & 0.720 & 0.760 \\
\bottomrule
\end{tabular}
\end{table}

Against integrated grade the panel outperforms arm-level SCNA in
\textbf{all twelve cells across both cohorts}, in both directions, roughly
tripling the rescue rate for over-called tumours.

\paragraph{The restricted cohort sharpens this to an absolute statement.}
Arm-level SCNA correctly reclassifies \textbf{none} of the integrated-grade
over-called tumours at any horizon (0/62, 0/34, 0/19), whereas the breakpoint
panel reclassifies 16\%, 24\% and 32\%. In a clean surgical population, arm-level
copy number carries no information at all about which elevated-integrated-grade
tumours will remain recurrence-free; positional breakpoint information carries
some. This is the single clearest demonstration in the analysis that the two
feature types are not measuring the same thing.

\paragraph{Against WHO grade the panel adds nothing} (Table~\ref{tab:discord_who}).
All three models sit at 0.33--0.55 over-called and 0.57--0.80 under-called, with
no consistent ordering. WHO grade's errors are already captured by arm-level copy
number; integrated grade's residual errors are not. Since integrated grade
already absorbs the arm-level signal (1p/22q loss, CDKN2A), what remains is
precisely the positional information the windows encode.

\note{Caution: restricted cells contain 5--73 tumours. The 3-year under-called
cell has $n=5$ and should not be interpreted.}

% =====================================================================
\section{The window list}
% =====================================================================

Selection frequency is the fraction of training folds that independently chose a
window; it separates a finding from a coincidence.

\begin{table}[H]
\centering
\caption{Windows selected in at least 40\% of training folds ($k=5$). Position
is arm-relative (0 = telomere, 1 = centromere).}
\label{tab:windows}
\begin{tabular}{llcccc}
\toprule
& & \multicolumn{3}{c}{Selection frequency} & \\
\cmidrule(lr){3-5}
Arm & Arm-relative & 3 yr & 5 yr & 7 yr & GRCh37 coordinates \\
\midrule
1p  & 0.90--0.95 & \textbf{1.00} & --- & 0.44 & chr1:112.50--118.75\,Mb \\
1p  & 0.95--1.00 & \textbf{0.96} & \textbf{0.88} & \textbf{0.92} & chr1:118.75--125.00\,Mb \\
6q  & 0.90--0.95 & \textbf{1.00} & \textbf{0.64} & --- & chr6:66.51--72.01\,Mb \\
6q  & 0.95--1.00 & 0.40 & --- & \textbf{0.68} & chr6:61.00--66.51\,Mb \\
19q & 0.50--0.55 & --- & \textbf{0.80} & \textbf{0.68} & chr19:41.18--42.81\,Mb \\
19q & 0.45--0.50 & --- & \textbf{0.72} & --- & chr19:42.81--44.45\,Mb \\
2p  & 0.75--0.80 & --- & \textbf{0.76} & --- & chr2:69.97--74.64\,Mb \\
4p  & 0.95--1.00 & 0.44 & --- & --- & chr4:47.88--50.40\,Mb \\
3p  & 0.50--0.55 & --- & --- & 0.40 & chr3:45.50--50.05\,Mb \\
\bottomrule
\end{tabular}
\end{table}

\paragraph{Positive control.}
The chromosome 1p 112--125\,Mb region is selected in 88--100\% of folds at every
horizon. This is the previously reported chr1p hotspot for this disease,
recovered here \emph{without} the frozen window list being available to the
selection procedure, and constitutes the strongest internal validation of the
pipeline.

\paragraph{Novel candidates.}
The chr19q 41--44.5\,Mb region (19q13) is selected in 68--80\% of folds at the
longer horizons and is the most interesting signal outside 1p. Chromosome 6q
loss is established meningioma biology. Windows below approximately 40\%
selection frequency should be treated as fold-specific noise.

% =====================================================================
\section{Model ladder and parsimony}
% =====================================================================

\begin{figure}[H]
\centering
\includegraphics[width=\textwidth]{figure6_full_ladder.png}
\caption{Full model ladder. Top row, genomic-only models; bottom row,
clinical-integrated models. Left, broad cohort; right, restricted. Clinical
covariates are age, $\log(1+\text{mitoses})$, gross-total resection, prior
radiotherapy and adjuvant radiotherapy.}
\label{fig:ladder}
\end{figure}

\begin{figure}[H]
\centering
\includegraphics[width=0.95\textwidth]{figure7_parsimony.png}
\caption{Discrimination against model size. The steep segment lies between 8 and
13 features; beyond 13 the curve is close to flat.}
\label{fig:parsimony}
\end{figure}

\begin{table}[H]
\centering
\caption{Model ladder by feature count, \textbf{broad cohort}. Bold marks the
best model at each horizon. Feature counts: integrated grade is a single
ordinal label; the panel is 8 binary window indicators; clinical adds 5
covariates; SCNA adds 47 arm-level calls; enriched adds 42 further breakpoint
features (amplitude, allelic, density).}
\label{tab:ladder}
\sisetup{table-format=1.3}
\begin{tabular}{lrSSS}
\toprule
Model & {Features} & {3 years} & {5 years} & {7 years} \\
\midrule
Integrated grade                       & 1   & 0.766 & 0.735 & 0.733 \\
Breakpoint panel alone                 & 8   & 0.702 & 0.735 & 0.752 \\
Clinical + panel                       & 13  & 0.836 & 0.830 & 0.819 \\
Clinical + SCNA + panel                & 60  & \bfseries 0.850 & 0.843 & 0.831 \\
Clinical + SCNA + enriched breakpoints & 102 & 0.844 & \bfseries 0.857 & \bfseries 0.849 \\
\bottomrule
\end{tabular}
\end{table}

\begin{table}[H]
\centering
\caption{Model ladder by feature count, \textbf{restricted cohort}
(primary, gross-total resection, no prior or adjuvant radiotherapy, not
radiation-induced, non-angiomatous). Estimates are transported from the broad
cohort without refitting; $n=387/248/155$ with $22/37/43$ recurrences.}
\label{tab:ladder_restricted}
\sisetup{table-format=1.3}
\begin{tabular}{lrSSS}
\toprule
Model & {Features} & {3 years} & {5 years} & {7 years} \\
\midrule
Integrated grade                       & 1   & \bfseries 0.827 & 0.782 & 0.752 \\
Breakpoint panel alone                 & 8   & 0.715 & 0.753 & 0.757 \\
Clinical + panel                       & 13  & 0.737 & 0.811 & 0.771 \\
Clinical + SCNA + panel                & 60  & 0.793 & \bfseries 0.843 & 0.802 \\
Clinical + SCNA + enriched breakpoints & 102 & 0.798 & 0.836 & \bfseries 0.822 \\
\bottomrule
\end{tabular}
\end{table}

Two observations follow in the broad cohort. \textbf{The largest single gain is
from 8 to 13 features} --- adding five routine clinical variables to the panel
is worth +0.09 to +0.13 AUC, far more than the 47 arm-level calls
(+0.012 to +0.014) or the 42 additional enriched breakpoint features
($-0.006$ to $+0.018$). \textbf{Thirteen features recover most of the achievable
performance}: 0.836 at 3 years against 0.844 for the 102-feature model.

\paragraph{The restricted cohort inverts this.}
Three differences are worth stating plainly
(Table~\ref{tab:ladder_restricted}). First, \textbf{integrated grade is the best
model at 3 years} (0.827), ahead of every genomic combination including the
102-feature model (0.798). Second, \textbf{arm-level SCNA is no longer
optional}: without it the 13-feature model reaches only 0.737 at 3 years,
well below integrated grade, and adding SCNA recovers +0.03 to +0.06 at every
horizon. Third, \textbf{the parsimonious panel matches or beats the enriched set
once SCNA is present} --- 0.843 versus 0.836 at 5 years --- which is what one
would expect when 102 features must be fitted against 37 events.

A likely explanation for the reversal is that both the windows and the clinical
covariates were selected in the broad cohort, where treatment and presentation
dominate recurrence; neither is tuned to a clean surgical setting. The 22 events
at 3 years also make that column the least precisely estimated in the report.

% =====================================================================
\section{External validation: a resolution failure}
% =====================================================================

The panel was applied to GSE189521 (110 tumours, EPIC methylation-derived copy
number), quarantined throughout and never used for discovery, thresholding or
model selection. Windows and coefficients were fixed from institutional data;
nothing was refitted.

\begin{table}[H]
\centering
\caption{External validation, GSE189521. Frozen transfer.}
\label{tab:external}
\sisetup{table-format=1.3}
\begin{tabular}{lrrSSS}
\toprule
Horizon & {$n$} & {Events} & {SCNA} & {SCNA + panel} & {Panel alone} \\
\midrule
3 years & 50 & 13 & 0.755 & 0.705 & 0.600 \\
5 years & 30 & 17 & 0.643 & 0.620 & 0.543 \\
7 years & 22 & 18 & 0.681 & 0.653 & 0.604 \\
\bottomrule
\end{tabular}
\end{table}

\textbf{The panel does not transfer; it degrades performance at every horizon.}
The mechanism was measured rather than inferred. Median breakpoints detected per
sample were 97 in the institutional cohort and 11 in the external cohort, and
mean occupancy of the eight panel windows fell from 27.8\% to 3.7\%, with two
windows firing in no external tumour at all. A marker defined at approximately
5\,Mb resolution cannot be evaluated on a platform detecting one-ninth as many
change points.

Notably, arm-level SCNA \emph{does} transfer (0.755 externally at 3 years,
against 0.763 for the previously published locked model). Coarse arm-level calls
are resolution-robust; fine positional windows are not. A fair external test
requires a same-resolution cohort.

% =====================================================================
\section{Limitations}
% =====================================================================

\begin{itemize}\itemsep2pt
\item Single institution, single platform, retrospective ascertainment of
      recurrence; surveillance intensity is not uniform.
\item Restricted-cohort estimates rest on 22, 37 and 43 events. The
      $+0.081$ increment at 5 years is the largest effect reported here and the
      least precisely estimated.
\item Windows are approximately 5--6\,Mb arm-relative deciles. They are regions
      of interest, not mapped rearrangement junctions.
\item Selection frequency measures stability across folds of a single cohort and
      is not external reproducibility.
\item The panel failed external validation on methylation-derived copy number
      and no cross-platform claim is made.
\item No multiplicity correction is applied to individual windows; they are
      reported by stability rather than by significance.
\item Research use only. No claim of clinical utility, treatment guidance or
      prospective validation is made.
\end{itemize}

% =====================================================================
\section{Reproducibility}\label{sec:repro}
% =====================================================================

All values were generated by scripts in \texttt{scripts/} and read directly from
the following files. Discrimination, specificity and reclassification analyses
use an identical model specification (elastic net, $\alpha=0.25$, $C=0.1$,
training-fold imputation); values shared between tables agree to
$<\!10^{-4}$.

\begin{table}[H]
\centering
\small
\begin{tabular}{lll}
\toprule
Content & Source file & Script \\
\midrule
Tables \ref{tab:disc}, \ref{tab:trend}, \ref{tab:ladder} & \texttt{full\_ladder.csv} & \texttt{full\_ladder.py} \\
Table \ref{tab:incr} & \texttt{clinical\_metrics\_discrimination.csv} & \texttt{hotlist\_clinical\_metrics.py} \\
Table \ref{tab:spec} & \texttt{clinical\_metrics\_operating.csv} & \texttt{hotlist\_clinical\_metrics.py} \\
Table \ref{tab:discord} & \texttt{grade\_discordance.csv} & \texttt{grade\_discordance.py} \\
Table \ref{tab:windows} & \texttt{hotlist\_windows.csv} & \texttt{discover\_hotlist.py} \\
Table \ref{tab:external} & \texttt{external\_validation.csv} & \texttt{external\_validate\_hotlist.py} \\
Comparator AUCs & \texttt{comparators.csv} & \texttt{compute\_comparators.py} \\
\bottomrule
\end{tabular}
\end{table}

Reproduction checks: WHO and integrated grade AUCs reproduce previously
published broad-cohort values to within 0.0005; restricted-cohort sizes
reproduce exactly (387/22, 248/37, 155/43); and arm-level SCNA in the restricted
cohort reproduces the published figure at 3 years (0.796 vs 0.779).

\end{document}
